% =============================================================================
% CHAPITRE 2 — ANALYSE DES BESOINS ET CAHIER DES CHARGES
% =============================================================================

\chapter{Analyse des besoins et cahier des charges}

% =============================================================================
% 2.1 INTRODUCTION
% =============================================================================

\section{Introduction}

Après avoir présenté dans le chapitre précédent le contexte général du
projet Droppy, sa problématique ainsi que les objectifs poursuivis, ce
chapitre est consacré à l'analyse des besoins et à la définition du cahier
des charges de la solution.

L'objectif de cette phase est de traduire le problème métier en exigences
concrètes permettant de guider la conception et la réalisation de la
plateforme. Il s'agit notamment d'identifier les utilisateurs concernés,
les sources de données, les fonctionnalités attendues ainsi que les
contraintes techniques et qualitatives auxquelles le système doit répondre.

L'analyse commence par une étude du domaine de la gestion des réseaux d'eau
et des problématiques liées aux fuites et aux comportements anormaux. Une
étude de l'existant permet ensuite de mettre en évidence les limites des
approches traditionnelles et de justifier la nécessité d'une supervision
plus automatisée et intelligente.

La vision fonctionnelle de Droppy est ensuite présentée afin de décrire la
chaîne globale allant de la collecte des données jusqu'à leur analyse et à
leur restitution à l'opérateur.

Enfin, les acteurs du système, les besoins fonctionnels, les besoins non
fonctionnels, les contraintes du projet ainsi que les critères de réussite
sont définis afin de constituer le cahier des charges de la solution.

% =============================================================================
% 2.2 COMPRÉHENSION DU DOMAINE MÉTIER
% =============================================================================

\section{Compréhension du domaine métier}

\subsection{Enjeux de la gestion des réseaux d'eau}

La gestion efficace d'un réseau de distribution d'eau nécessite une
surveillance régulière des différents équipements et des flux de
consommation. Le gestionnaire doit être capable de suivre l'évolution des
débits, des consommations ainsi que l'état des équipements afin
d'identifier rapidement toute situation inhabituelle.

Dans un réseau instrumenté, les capteurs permettent de transformer les
équipements physiques en sources de données exploitables par une plateforme
logicielle. La surveillance ne repose alors plus uniquement sur des
inspections ponctuelles, mais peut s'appuyer sur une observation continue
du comportement du réseau.

Cette évolution présente un intérêt particulier dans le contexte de la
détection des fuites. Une fuite non détectée peut entraîner une perte
importante d'eau et provoquer des coûts supplémentaires pour l'exploitant.
La rapidité de détection constitue donc un élément essentiel de la
supervision.

\subsection{Typologie des anomalies}

Dans le cadre du projet Droppy, une anomalie correspond à un comportement
qui s'écarte du fonctionnement habituellement observé.

Plusieurs situations peuvent notamment être considérées :

\begin{itemize}

    \item une augmentation inhabituelle du débit ;

    \item une consommation anormalement élevée pendant une période de faible
    activité ;

    \item une consommation nocturne inhabituelle ;

    \item une incohérence entre l'état d'une vanne et le débit observé ;

    \item une évolution inhabituelle d'une série temporelle ;

    \item un comportement globalement atypique détecté automatiquement par
    un modèle d'apprentissage automatique.

\end{itemize}

Il est important de souligner qu'une anomalie ne correspond pas
nécessairement à une fuite. Elle constitue plutôt un signal nécessitant
une analyse complémentaire. La plateforme doit donc fournir suffisamment
d'informations pour permettre à l'opérateur d'interpréter le comportement
détecté.

\subsection{Données disponibles et télémétrie IoT}

Les capteurs IoT constituent la principale source de données de Droppy.

Les informations exploitées peuvent notamment comprendre :

\begin{itemize}

    \item le débit d'eau ;

    \item la consommation ;

    \item l'état d'une vanne ;

    \item l'identifiant du dispositif ;

    \item l'horodatage associé à la mesure.

\end{itemize}

Chaque équipement doit pouvoir être identifié de manière unique afin de
maintenir la correspondance entre les mesures reçues et leur source.

La dimension temporelle des données constitue également un élément
important. L'analyse d'une mesure isolée étant généralement insuffisante,
le système doit pouvoir exploiter l'évolution des données au cours du
temps.

\subsection{Enjeux de la détection précoce}

La détection précoce représente un enjeu majeur pour la surveillance des
réseaux d'eau.

Une anomalie détectée rapidement permet à l'opérateur d'analyser la
situation et, lorsque cela est nécessaire, de déclencher une intervention
avant que le problème ne provoque des pertes importantes.

L'objectif de Droppy est donc de réduire autant que possible le délai entre
l'apparition d'un comportement inhabituel et son identification par
l'utilisateur.

Cette exigence justifie la mise en place d'une chaîne automatisée combinant
collecte des données, analyse, détection et notification.

% =============================================================================
% 2.3 ÉTUDE DE L'EXISTANT
% =============================================================================

\section{Étude de l'existant}

\subsection{Surveillance traditionnelle}

Les approches traditionnelles de surveillance des réseaux d'eau reposent
principalement sur les inspections périodiques, les relevés manuels et
l'utilisation de seuils définis à l'avance.

Dans ce contexte, l'opérateur consulte les mesures disponibles et recherche
des valeurs pouvant indiquer un dysfonctionnement.

Cette approche reste pertinente pour certains scénarios, mais elle devient
plus difficile à maintenir lorsque le nombre d'équipements et la fréquence
des mesures augmentent.

\subsection{Approches basées sur les seuils}

Une méthode courante consiste à définir un seuil à partir duquel une alerte
est déclenchée.

Par exemple, une règle peut considérer qu'un débit dépassant une certaine
valeur constitue une situation potentiellement anormale.

Cependant, le comportement normal d'un équipement peut varier selon
l'heure, le jour, la période d'utilisation ou les caractéristiques du site.
Un seuil fixe ne permet donc pas toujours de représenter correctement
l'ensemble des situations possibles.

\subsection{Limites identifiées}

Plusieurs limites peuvent ainsi être identifiées.

Premièrement, les inspections manuelles ne garantissent pas une détection
immédiate des anomalies.

Deuxièmement, les seuils statiques peuvent générer des faux positifs lorsque
la variation observée correspond à un fonctionnement normal.

Troisièmement, certaines anomalies peuvent rester invisibles lorsqu'elles
ne dépassent pas les seuils prédéfinis.

Enfin, l'analyse manuelle devient difficile lorsque plusieurs capteurs
produisent continuellement des données.

Ces limites montrent qu'une approche automatisée est nécessaire afin
d'exploiter efficacement les volumes de données disponibles.

\subsection{Besoin d'une supervision intelligente}

Une supervision intelligente doit être capable de combiner plusieurs
sources d'information et plusieurs mécanismes d'analyse.

Dans Droppy, cette approche repose sur une combinaison entre des règles
métier et des techniques d'apprentissage automatique.

Les règles métier permettent d'identifier des situations connues, tandis
que les modèles d'apprentissage automatique permettent de rechercher des
comportements atypiques dans les données.

Cette complémentarité permet d'obtenir une analyse plus riche qu'une
approche reposant uniquement sur des seuils fixes.

% =============================================================================
% 2.4 VISION FONCTIONNELLE DE DROPPY
% =============================================================================

\section{Vision fonctionnelle de Droppy}

\subsection{Vision globale}

Droppy est conçu comme une plateforme de supervision intelligente permettant
de transformer les données brutes provenant des équipements IoT en
informations exploitables par les opérateurs.

La solution est organisée autour de trois composants logiciels principaux :

\begin{itemize}

    \item \textbf{Angular 19}, utilisé pour développer le tableau de bord
    de supervision ;

    \item \textbf{Spring Boot 3}, utilisé comme service backend central pour
    l'intégration, l'orchestration et la gestion des événements ;

    \item \textbf{FastAPI et Python}, utilisés pour les traitements de
    données et les mécanismes d'intelligence artificielle.

\end{itemize}

Le dépôt du projet confirme cette organisation en distinguant notamment
le frontend \texttt{dashboard-angular}, le service backend
\texttt{core-service} et le moteur d'intelligence artificielle situé dans
le répertoire \texttt{src}. :contentReference[oaicite:3]{index=3}

\subsection{Chaîne globale de traitement}

Le fonctionnement fonctionnel de Droppy peut être représenté comme une
chaîne composée des étapes suivantes :

\begin{enumerate}

    \item collecte des données provenant des équipements IoT ;

    \item ingestion des données par le système ;

    \item nettoyage et préparation des données ;

    \item application des règles métier ;

    \item analyse par les modèles d'apprentissage automatique ;

    \item génération des résultats de détection ou de prévision ;

    \item enregistrement des informations importantes ;

    \item transmission des événements au tableau de bord ;

    \item consultation et interprétation des résultats par l'opérateur.

\end{enumerate}

Cette chaîne permet de transformer progressivement une donnée brute en une
information pouvant être utilisée pour la supervision et la prise de
décision.

\subsection{Valeur ajoutée de la solution}

La valeur ajoutée de Droppy réside principalement dans l'intégration de
plusieurs mécanismes au sein d'une même plateforme.

La solution ne se limite pas à afficher les mesures provenant des capteurs.
Elle permet également de les nettoyer, de les analyser et d'identifier
automatiquement certains comportements inhabituels.

L'utilisation d'un modèle tel que \textbf{Isolation Forest} permet notamment
de rechercher des observations atypiques sans dépendre exclusivement de
seuils prédéfinis.

La plateforme intègre également des mécanismes de prévision ainsi que des
notifications permettant de transmettre rapidement les événements
importants à l'utilisateur.

Le dépôt du projet indique notamment l'existence des fonctionnalités de
réentraînement des modèles, de prédiction combinée et de prévision de
consommation. :contentReference[oaicite:4]{index=4}

\subsection{Résultats attendus}

À partir des données collectées, le système doit permettre d'obtenir
plusieurs types de résultats :

\begin{itemize}

    \item identification de comportements atypiques ;

    \item détection de situations pouvant correspondre à des fuites ;

    \item production de prévisions de consommation ;

    \item génération d'événements d'alerte ;

    \item visualisation des résultats dans le tableau de bord ;

    \item conservation des informations nécessaires à l'analyse historique.

\end{itemize}

L'objectif final est de fournir à l'opérateur une information synthétique
et exploitable plutôt qu'un simple ensemble de mesures brutes.

% =============================================================================
% 2.5 IDENTIFICATION DES ACTEURS
% =============================================================================

\section{Identification des acteurs}

L'identification des acteurs permet de déterminer les différentes entités
participant au fonctionnement de la solution et leurs responsabilités.

Contrairement à une application classique, Droppy implique à la fois des
utilisateurs humains, des équipements physiques et des services logiciels.

\subsection{Opérateur / gestionnaire}

L'opérateur constitue l'utilisateur principal de la plateforme.

Il utilise le tableau de bord pour :

\begin{itemize}

    \item consulter les équipements ;

    \item visualiser les données de télémétrie ;

    \item suivre les anomalies détectées ;

    \item consulter les historiques ;

    \item recevoir les notifications ;

    \item analyser les tendances de consommation.

\end{itemize}

Son rôle est principalement orienté vers la supervision et l'interprétation
des résultats fournis par le système.

\subsection{Capteurs et équipements IoT}

Les équipements IoT représentent les sources physiques des données.

Ils transmettent les mesures nécessaires au fonctionnement de la chaîne
d'analyse. Ces données peuvent notamment concerner le débit, la
consommation et l'état des vannes.

Les équipements doivent être identifiables afin de permettre au système
d'associer chaque mesure au dispositif correspondant.

\subsection{Backend applicatif}

Le backend Spring Boot joue le rôle de service central d'intégration.

Il assure notamment :

\begin{itemize}

    \item la réception et l'intégration des données ;

    \item l'orchestration des traitements ;

    \item la communication avec le service d'intelligence artificielle ;

    \item la persistance des informations ;

    \item le routage des événements vers le tableau de bord.

\end{itemize}

Le dépôt indique notamment l'utilisation de Spring Boot 3, de mécanismes de
planification et de WebSockets pour les notifications en temps réel. :contentReference[oaicite:5]{index=5}

\subsection{Service d'intelligence artificielle}

Le microservice FastAPI constitue le moteur d'analyse de Droppy.

Il prend en charge les traitements liés aux données et à l'intelligence
artificielle, notamment :

\begin{itemize}

    \item le nettoyage des données ;

    \item le prétraitement ;

    \item l'application de certaines règles métier ;

    \item la détection des anomalies ;

    \item la prévision de consommation ;

    \item le réentraînement des modèles.

\end{itemize}

Le projet expose notamment des opérations de réentraînement, de prédiction
et de prévision pour un équipement donné. :contentReference[oaicite:6]{index=6}

\subsection{Interface de supervision}

L'interface Angular constitue le point d'accès principal de l'opérateur.

Elle permet de centraliser les résultats produits par les différents
services et de les présenter sous une forme visuelle.

Le tableau de bord permet notamment d'afficher les anomalies, les données
historiques et les graphiques de télémétrie tout en recevant les
notifications en temps réel. :contentReference[oaicite:7]{index=7}

% =============================================================================
% 2.6 BESOINS FONCTIONNELS
% =============================================================================

\section{Besoins fonctionnels}

Les besoins fonctionnels définissent les services que la plateforme doit
fournir pour répondre aux objectifs identifiés.

\subsection{Gestion des équipements}

\textbf{BF01 -- Gestion des équipements}

Le système doit permettre d'identifier et de consulter les équipements
surveillés.

Chaque équipement doit pouvoir être associé à ses différentes mesures et
à son historique afin de faciliter le suivi individuel.

\subsection{Ingestion des télémétries}

\textbf{BF02 -- Collecte des données}

Le système doit permettre de recevoir et d'intégrer les données produites
par les équipements IoT.

Les données doivent être correctement associées à leur équipement et à
leur instant de mesure.

\subsection{Nettoyage et prétraitement}

\textbf{BF03 -- Prétraitement}

Le système doit être capable de nettoyer et de préparer les données avant
leur analyse.

Cette étape doit notamment permettre de traiter les valeurs manquantes,
les incohérences et les transformations nécessaires à l'exploitation des
séries temporelles.

\subsection{Détection des anomalies}

\textbf{BF04 -- Détection des anomalies}

La plateforme doit identifier automatiquement les comportements atypiques
présents dans les données de télémétrie.

Un modèle de type \textbf{Isolation Forest} est utilisé pour contribuer à
cette détection.

\subsection{Détection des fuites}

\textbf{BF05 -- Détection des fuites}

Le système doit permettre d'identifier les situations pouvant correspondre
à une fuite.

La détection peut s'appuyer sur plusieurs indicateurs tels que le débit,
la consommation, l'état des vannes et les comportements observés pendant
les périodes de faible activité.

\subsection{Prévision de consommation}

\textbf{BF06 -- Prévision}

La plateforme doit permettre de produire des prévisions de consommation
afin de comparer les valeurs observées avec les tendances attendues.

Cette fonctionnalité peut également contribuer à identifier des écarts
significatifs dans le comportement d'un équipement.

\subsection{Gestion des alertes}

\textbf{BF07 -- Alertes}

Lorsqu'une anomalie importante est identifiée, le système doit pouvoir
générer un événement destiné à l'opérateur.

Les événements critiques doivent pouvoir être transmis rapidement au
tableau de bord.

\subsection{Supervision en temps réel}

\textbf{BF08 -- Supervision}

Le système doit fournir une interface permettant à l'opérateur de suivre
les informations importantes relatives au réseau.

La plateforme doit notamment permettre de consulter les anomalies, les
graphiques, les équipements et les notifications.

\subsection{Historisation}

\textbf{BF09 -- Historique}

Le système doit conserver les informations nécessaires au suivi historique
des données et des anomalies.

L'historique permet notamment de comparer le comportement actuel d'un
équipement avec ses observations précédentes.

\subsection{Réentraînement des modèles}

\textbf{BF10 -- Gestion des modèles}

Le système doit permettre de réentraîner les modèles d'intelligence
artificielle pour un équipement donné lorsque cela est nécessaire.

Cette fonctionnalité permet d'adapter les modèles aux nouvelles données
disponibles.

% =============================================================================
% 2.7 BESOINS NON FONCTIONNELS
% =============================================================================

\section{Besoins non fonctionnels}

\subsection{Performance}

\textbf{BNF01 -- Performance}

La plateforme doit traiter les données avec un temps de réponse acceptable
afin de permettre une utilisation fluide du tableau de bord.

\subsection{Fiabilité}

\textbf{BNF02 -- Fiabilité}

Les données doivent être traitées de manière cohérente et les mécanismes
d'analyse doivent limiter les erreurs liées aux données incomplètes ou
incohérentes.

\subsection{Temps réel}

\textbf{BNF03 -- Temps réel}

Les événements importants doivent pouvoir être transmis rapidement à
l'interface de supervision.

L'utilisation des WebSockets permet au backend de transmettre des
notifications sans nécessiter une actualisation permanente de la page.

\subsection{Sécurité}

\textbf{BNF04 -- Sécurité}

Les communications entre les différents composants doivent être protégées
et les accès aux fonctionnalités sensibles doivent être contrôlés.

\subsection{Maintenabilité}

\textbf{BNF05 -- Maintenabilité}

La séparation entre frontend, backend et service d'intelligence artificielle
doit faciliter la maintenance et l'évolution du système.

\subsection{Évolutivité}

\textbf{BNF06 -- Évolutivité}

La solution doit pouvoir intégrer de nouveaux équipements, de nouvelles
méthodes d'analyse et de nouvelles fonctionnalités sans nécessiter une
refonte complète de l'architecture.

\subsection{Ergonomie}

\textbf{BNF07 -- Ergonomie}

Le tableau de bord doit présenter les informations de manière claire et
compréhensible afin de permettre à l'opérateur d'identifier rapidement
les situations nécessitant son attention.

% =============================================================================
% 2.8 CONTRAINTES DU PROJET
% =============================================================================

\section{Contraintes du projet}

La réalisation de Droppy doit respecter plusieurs contraintes techniques
et organisationnelles.

\subsection{Hétérogénéité technologique}

La solution repose sur plusieurs technologies et environnements :

\begin{itemize}

    \item Python et FastAPI pour le service d'intelligence artificielle ;

    \item Java et Spring Boot pour le backend ;

    \item TypeScript et Angular 19 pour le frontend ;

    \item outils de gestion de données et de modèles pour le pipeline IA.

\end{itemize}

Cette diversité nécessite une définition claire des interfaces entre les
différents composants.

\subsection{Nature temporelle des données}

Les données de télémétrie sont associées à des instants de mesure. Leur
ordre temporel doit donc être conservé afin de permettre l'analyse des
évolutions et des tendances.

\subsection{Communication inter-services}

Le backend doit assurer une communication fiable avec le service
d'intelligence artificielle et avec le frontend.

Le dépôt du projet indique notamment l'utilisation de HTTP/REST entre les
services ainsi que de WebSockets pour les notifications temps réel. :contentReference[oaicite:8]{index=8}

\subsection{Gestion de la persistance}

Les informations nécessaires au suivi des équipements, des événements et
des analyses doivent être correctement persistées.

Le projet utilise notamment Liquibase pour la gestion et le versionnement
des évolutions du schéma de base de données. :contentReference[oaicite:9]{index=9}

\subsection{Validation de la solution}

La solution doit être accompagnée de tests permettant de vérifier les
différents composants.

Le dépôt comprend notamment des tests dédiés au pipeline Python et aux
fonctionnalités de l'API, ainsi que des procédures de test pour les
différents composants de l'application. :contentReference[oaicite:10]{index=10}

% =============================================================================
% 2.9 CAHIER DES CHARGES FONCTIONNEL
% =============================================================================

\section{Cahier des charges fonctionnel}

Le cahier des charges fonctionnel synthétise les principales exigences
identifiées au cours de l'analyse.

\subsection{Tableau des exigences}

\begin{table}[H]
\centering
\caption{Synthèse des exigences fonctionnelles de Droppy}
\label{tab:exigences-droppy}

\renewcommand{\arraystretch}{1.35}

\begin{tabularx}{\textwidth}{
    >{\bfseries\centering\arraybackslash}p{1.5cm}
    >{\raggedright\arraybackslash}p{3.1cm}
    >{\raggedright\arraybackslash}X
    >{\centering\arraybackslash}p{1.8cm}
}

\toprule

ID & Fonctionnalité & Description & Priorité \\

\midrule

BF01 &
Gestion des équipements &
Identifier et consulter les équipements surveillés. &
Haute \\

BF02 &
Collecte des données &
Recevoir et intégrer les télémétries provenant des capteurs. &
Haute \\

BF03 &
Prétraitement &
Nettoyer et préparer les données avant leur analyse. &
Haute \\

BF04 &
Détection des anomalies &
Identifier automatiquement les comportements atypiques. &
Haute \\

BF05 &
Détection des fuites &
Identifier les situations pouvant correspondre à une fuite. &
Haute \\

BF06 &
Prévision &
Analyser les tendances et prévoir la consommation. &
Moyenne \\

BF07 &
Alertes &
Informer l'opérateur lorsqu'une anomalie importante est détectée. &
Haute \\

BF08 &
Supervision &
Visualiser les données, équipements et anomalies. &
Haute \\

BF09 &
Historique &
Permettre la consultation des données historiques. &
Moyenne \\

BF10 &
Gestion des modèles &
Permettre le réentraînement et l'exploitation des modèles. &
Moyenne \\

\bottomrule

\end{tabularx}

\end{table}

\subsection{Priorisation des exigences}

Les exigences ont été classées selon leur importance pour le fonctionnement
global de la solution.

Les fonctionnalités de priorité \textbf{haute} constituent le noyau de
Droppy. Elles comprennent notamment la collecte des données, le
prétraitement, la détection des anomalies, la détection des fuites, les
alertes et la supervision.

Les fonctionnalités de priorité \textbf{moyenne} viennent compléter le
système et permettent d'améliorer les capacités d'analyse et d'exploitation
des données.

Cette priorisation permet de concentrer les efforts de développement sur
les fonctionnalités essentielles avant d'intégrer les fonctionnalités
complémentaires.

% =============================================================================
% 2.10 CRITÈRES DE RÉUSSITE
% =============================================================================

\section{Critères de réussite}

La réussite de la plateforme Droppy est évaluée selon plusieurs dimensions.

\subsection{Critères fonctionnels}

Le système doit être capable :

\begin{itemize}

    \item de recevoir correctement les données de télémétrie ;

    \item de traiter et préparer les données ;

    \item d'identifier les comportements atypiques ;

    \item de détecter les situations pouvant correspondre à des fuites ;

    \item de produire des prévisions ;

    \item de générer et transmettre les alertes ;

    \item d'afficher les résultats dans le tableau de bord.

\end{itemize}

\subsection{Critères techniques}

La solution doit également respecter plusieurs critères techniques :

\begin{itemize}

    \item communication correcte entre les différents services ;

    \item fonctionnement des échanges temps réel ;

    \item persistance fiable des informations ;

    \item possibilité de réentraîner les modèles ;

    \item modularité des composants ;

    \item présence de mécanismes de test et de validation.

\end{itemize}

\subsection{Critères liés à l'utilisateur}

Enfin, le système doit permettre à l'opérateur :

\begin{itemize}

    \item d'accéder rapidement aux informations importantes ;

    \item d'identifier facilement les anomalies ;

    \item de consulter l'historique ;

    \item de comprendre les résultats produits ;

    \item de réagir rapidement lorsqu'une alerte critique est générée.

\end{itemize}

% =============================================================================
% 2.11 CONCLUSION
% =============================================================================

\section{Conclusion}

Ce chapitre a permis de transformer la problématique générale de la
détection des fuites d'eau en un ensemble d'exigences fonctionnelles et
non fonctionnelles concrètes.

L'analyse du domaine métier a montré l'importance de disposer d'une
surveillance continue capable d'exploiter les données issues des
équipements IoT. L'étude de l'existant a également mis en évidence les
limites des inspections périodiques et des approches reposant uniquement
sur des seuils fixes.

La vision fonctionnelle de Droppy repose ainsi sur une chaîne complète
allant de la collecte des télémétries jusqu'à leur analyse et leur
restitution dans une interface de supervision. Cette chaîne combine
traitement des données, règles métier, apprentissage automatique,
prévision et notifications.

Les acteurs du système ainsi que les besoins fonctionnels, les besoins non
fonctionnels et les contraintes ont ensuite été identifiés. Le cahier des
charges obtenu constitue la base de la phase de conception.

Le chapitre suivant sera consacré à la \textbf{conception de la solution}.
Il présentera l'architecture globale de Droppy, les interactions entre ses
différents composants, les flux de données ainsi que les modèles permettant
de formaliser la conception du système.